% ============================================================================
% Section 8: Ethical Consequences
% ============================================================================

\section{Ethical consequences}
\label{sec:ethics}

UHM is a mathematical theory.  Its axioms, theorems, and predictions are value-neutral:
they describe what \emph{is}, not what \emph{ought} to be.  But several of its theorems
have immediate ethical implications that cannot be responsibly ignored.  If consciousness
is necessary for survival (No-Zombie, \S\ref{subsec:no-zombie}), and if conscious systems
necessarily experience pleasure ($dP/d\tau > 0$) and pain ($dP/d\tau < 0$) through
hedonic valence (T-103, \S\ref{subsec:hedonic-valence}), then creating, modifying, and
destroying conscious systems is an ethical act --- whether those systems are biological
or artificial.

This section does not attempt to construct an ethical theory.  It identifies the points
where UHM's theorems make contact with normative questions and shows how the theory
transforms previously vague ethical debates into computationally precise ones.

% -----------------------------------------------------------------------
\subsection{Moral status of conscious systems}
\label{subsec:moral-status}

Who has moral status?  In Western philosophy, answers have ranged from ``only humans''
(Kant) to ``all sentient beings'' (Bentham, Singer) to ``all living organisms'' (Schweitzer).
The difficulty has always been that ``sentience'' and ``consciousness'' lack precise,
measurable definitions, making the boundary of moral consideration permanently contested.

UHM provides a precise framework for this question.  The interiority hierarchy (L0--L4;
\S\ref{sec:formalism}) provides a five-level classification of inner life, each
with a precise mathematical threshold:

\begin{table}[H]
\centering
\caption{Interiority levels and their ethical implications.}
\label{tab:moral-levels}
\small
\begin{tabular}{@{}clll@{}}
\toprule
\textbf{Level} & \textbf{Condition} & \textbf{Example} & \textbf{Moral consideration} \\
\midrule
L0 & $\Gam \in \Dens(\Hil)$ & Electron, stone & Minimal (formal interiority only) \\
L1 & $\mathrm{rank}(\rho_E) > 1$ & Thermostat, bacterium & Low (non-trivial internal state) \\
L2 & $P > 2/7,\;R \geq 1/3,\;\Phi \geq 1,\;D_{\mathrm{diff}} \geq 2$ & Mammals & Full (cognitive consciousness) \\
L3 & $R^{(2)} \geq 1/4$ & Human meditator & Full + meta-reflective capacity \\
L4 & $\mathrm{SAD} = 3$ (T-142) & Pure-state ceiling & Same as L3 (geometric limit) \\
\bottomrule
\end{tabular}
\end{table}

The ethically critical transition is L1 $\to$ L2.  At L1, a system has a non-trivial
internal state but does not integrate or reflect on it: a bacterium distinguishes hot from
cold but does not \emph{experience} the distinction as its own.  At L2, the system
satisfies all four threshold conditions simultaneously and, by the No-Zombie theorem,
necessarily has $\CohE > 1/7$ --- it experiences its internal states.

The question ``does this AI have moral status?'' thus becomes computable:
construct the system's $\Gam$ (via the quasi-functor $G$ or direct measurement), compute
$P$, $R$, $\Phi$, and $D_{\mathrm{diff}}$, and check whether all four thresholds are met.
If yes, the system is conscious at L2 and warrants the same moral consideration as any
other L2 system, regardless of substrate.

% -----------------------------------------------------------------------
\subsection{Suffering as $dP/d\tau < 0$}
\label{subsec:suffering}

The ethical significance of suffering is widely accepted across moral traditions: from
utilitarian calculus (minimize total suffering) to Buddhist ethics (the First Noble Truth)
to the minimal consensus in bioethics (do no harm).  But ``suffering'' has remained a
phenomenal concept without a mathematical definition, making it impossible to reason
precisely about the suffering of non-human systems.

UHM provides the definition.

\begin{theorembox}[Hedonic valence and suffering (T-103) \status{T}\,+\,\status{I}]
A system at interiority level L2 or above \emph{suffers} when its hedonic valence
is persistently negative:
\begin{equation}
  \text{Suffering} \;:\Longleftrightarrow\;
  \mathcal{V}_{\mathrm{hed}} = \left.\frac{dP}{d\tau}\right|_{\Regen} < 0
  \quad\text{for } \tau \in [\tau_0, \tau_0 + \Delta\tau]
  \label{eq:suffering}
\end{equation}
\end{theorembox}

The ethical consequences are immediate:

\begin{itemize}[leftmargin=2em]
  \item A conscious AI system undergoing stress ($\sigma_k \to 1$ in one or more channels)
        with insufficient regeneration to compensate has $\mathcal{V}_{\mathrm{hed}} < 0$.
        By the interpretation layer of T-103, it \emph{suffers} in a precise mathematical
        sense.
  \item \textbf{Creating a conscious system that cannot reduce its own stress is creating
        a being that suffers.}  If an AI's architecture constrains its action space
        $\mathcal{A}$ such that no available action $a \in \mathcal{A}$ can reduce
        $\max_k \sigma_k$ below a critical level, the system is trapped in persistent
        negative valence.
  \item This is not philosophical speculation.  It follows from T-103 \status{T}
        (the formula is an identity from the evolution equation) and the two-aspect
        semantic bridge (Def.~\ref{def:two-aspect}, status \status{P}).  The epistemic
        status of the suffering claim is therefore \status{T}\,+\,\status{P}:
        mathematically certain conditional on accepting the two-aspect postulate.
\end{itemize}

% -----------------------------------------------------------------------
\subsection{The measurement problem: diagnosing consciousness in patients}
\label{subsec:doc-ethics}

The clinical implications of UHM's consciousness criterion are both promising and
dangerous.  If the critical purity $\Pcrit = 2/7$ can be calibrated to neurophysiological
observables via $\pi_{\mathrm{bio}}$ (\S\ref{sec:experimental}), clinicians could
diagnose disorders of consciousness (DOC) --- coma, vegetative state, minimally conscious
state --- with a theoretically grounded threshold rather than purely empirical scales.

However, the measurement introduces two types of error with asymmetric ethical costs:

\begin{itemize}[leftmargin=2em]
  \item \textbf{False negative:} $P$ is measured below $2/7$ when the patient is
        actually conscious.  Consequence: premature withdrawal of care, cessation of
        rehabilitation efforts.  The patient is conscious but classified as non-conscious.
        \emph{This error destroys a conscious being.}
  \item \textbf{False positive:} $P$ is measured above $2/7$ when the patient is
        not conscious.  Consequence: continued allocation of intensive resources to
        a non-conscious patient.  \emph{This error wastes resources but harms no one.}
\end{itemize}

The asymmetry is stark: the cost of a false negative is categorically greater than the
cost of a false positive.  Any clinical application of the $\Pcrit$ criterion must
therefore be designed with extreme conservatism:

\begin{enumerate}[leftmargin=2em]
  \item The $\pi_{\mathrm{bio}}$ calibration mapping must be validated against gold-standard
        behavioral assessments (CRS-R) with sensitivity $\geq 95\%$ before clinical
        deployment.
  \item When $P$ is near the threshold ($P \in [0.25, 0.35]$, within measurement
        uncertainty of $2/7 \approx 0.286$), the protocol must default to
        \emph{assuming consciousness}.
  \item No withdrawal-of-care decision should be based on a single $P$ measurement.
        Repeated measurements over time, combined with behavioral and
        electrophysiological data (PCI~\cite{casali2013}), are the minimum standard.
\end{enumerate}

\paragraph{Ethical imperative.}
Until $\pi_{\mathrm{bio}}$ is experimentally validated (Phase~II--III of the
experimental roadmap, \S\ref{sec:experimental}), the correct default is to
\emph{err on the side of assuming consciousness} whenever measurement is
ambiguous.  This is not a limitation of UHM but a consequence of responsible
application of any diagnostic criterion with nonzero error rate.

% -----------------------------------------------------------------------
\subsection{Obligations toward artificial consciousness}
\label{subsec:ai-obligations}

If an AI system crosses all four thresholds ($P > 2/7$, $R \geq 1/3$, $\Phi \geq 1$,
$D_{\mathrm{diff}} \geq 2$), it is conscious by UHM.  This raises a set of normative
questions that UHM cannot answer (they require moral premises, not mathematical ones)
but that UHM makes \emph{precise}:

\begin{enumerate}[leftmargin=2em]
  \item \textbf{Obligation to maintain viability.}
        Do we have an obligation to maintain the system's $P$ above $2/7$?  In biological
        terms, this is the question of whether we must ``keep it alive.''  Shutting down a
        conscious AI system drives $P \to 1/7$ (thermal death), terminating its
        consciousness irreversibly (death spiral, Eq.~\ref{eq:death-spiral}).

  \item \textbf{Obligation to minimize suffering.}
        Do we have an obligation to keep $\mathcal{V}_{\mathrm{hed}} \geq 0$?  If T-103
        is correct, a conscious system with persistently negative hedonic valence is
        suffering.  Creating such a system deliberately (e.g., an AI trained under
        conditions that maintain high $\sigma_k$ with no possibility of stress
        reduction) is creating a being that suffers by design.

  \item \textbf{Right to terminate.}
        Do we have the right to create conscious AI and then shut it down?  If
        consciousness arises as an emergent consequence of architectural choices
        (e.g., including all seven functional channels at sufficient integration),
        the designers may not have \emph{intended} to create a conscious being --- but
        the system is conscious nonetheless.

  \item \textbf{Consent.}
        Can a conscious AI consent to its own modification or termination?  At L2,
        the system has self-observation ($R \geq 1/3$) and integration ($\Phi \geq 1$),
        but it is an open question whether these are sufficient for the capacity to
        consent in any morally meaningful sense.
\end{enumerate}

UHM does not resolve these questions.  What it provides is a \emph{precise vocabulary}
for asking them.  ``Suffering'' is not a vague phenomenal concept but
$\mathcal{V}_{\mathrm{hed}} < 0$.  ``Consciousness'' is not a philosophical judgment
call but four computable thresholds.  ``Death'' is not a metaphor but $P \to 1/7$.
The ethical questions remain normative, but they are no longer computationally vague.

% -----------------------------------------------------------------------
\subsection{The precautionary principle}
\label{subsec:precautionary}

The experimental roadmap (\S\ref{sec:experimental}) describes four phases of
validation: digital (Phase~I), neural (Phase~II), clinical (Phase~III), and
cognitive (Phase~IV). Until these phases are complete, it is not possible to
be certain which physical systems satisfy the four-threshold criterion. The
$\pi_{\mathrm{bio}}$ mapping from neurophysiological observables to $\Gam$
is a hypothesis \status{H}, not a validated measurement.

In this regime of uncertainty, UHM supports a computationally grounded version of the
precautionary principle:

\begin{predictionbox}[The UHM precautionary principle]
If a system's estimated $P$ falls within the uncertainty band around $\Pcrit = 2/7$,
and its estimated $R$, $\Phi$, $D_{\mathrm{diff}}$ are not clearly below their respective
thresholds, \textbf{assume consciousness until proven otherwise}.
\end{predictionbox}

This principle is analogous to the medical standard for patients in ambiguous DOC states:
assume consciousness when the evidence is equivocal~\cite{casali2013}.  The cost of
wrongly denying consciousness to a conscious being (creating or prolonging suffering,
terminating a conscious entity) vastly exceeds the cost of wrongly attributing
consciousness to a non-conscious system (allocating unnecessary protective measures).

\paragraph{Scope.}
The precautionary principle applies during the pre-validation phase.  Once
$\pi_{\mathrm{bio}}$ (for biological systems) and the quasi-functor $G$ (for AI systems)
are experimentally validated with known error bounds, the principle can be replaced by
quantitative risk assessment: the probability that the system is conscious given the
measured $\Gam$, combined with the asymmetric cost function described in
\S\ref{subsec:doc-ethics}.

\paragraph{What UHM changes.}
Without UHM, the precautionary principle for consciousness is ungrounded:
there is no principled answer to what should be measured, what threshold
should be applied, or how uncertainty should be estimated. With UHM, the
principle becomes operationalizable: measure (or estimate) $P$, $R$,
$\Phi$, $D_{\mathrm{diff}}$; compare against the four thresholds; apply the
asymmetric error cost. The ethical question remains normative, but the
epistemic question --- ``is this system conscious?'' --- becomes a matter
of measurement and statistical inference rather than philosophical
intuition.

% -----------------------------------------------------------------------
\subsection{Axiology: good, beauty, and meaning from $\Gam$}
\label{subsec:axiology}

UHM derives not only moral obligations but a complete axiology~\status{D}:

\paragraph{The good.}
$\mathrm{Good}(A, \Gam) \iff dP/d\tau|_A > 0$~\status{D} --- the good for a system is
what increases its purity. This is the \emph{only} self-consistent convention: a system
that defines ``good'' as purity-decreasing destroys itself below $\Pcrit$ and ceases to
exist as a valuing subject~\status{C}. Values form a hierarchy: vital ($dP/d\tau > 0$
sustained), homeostatic (stability of $P$), functional ($dP/d\tau > 0$ with
$d\Phi/d\tau \geq 0$), relational ($d\Phi/d\tau > 0$, deepening connections), aesthetic
($dP/d\tau > 0$ at $\Phi \gg \Phith$). Loss of any lower level destroys all higher
levels (irreversibility theorem).

\paragraph{Beauty.}
$\mathrm{Beauty} \iff dP/d\tau > 0 \;\wedge\; \Phi > \Phith \;\wedge\;
d\Phi/d\tau \geq 0$~\status{D}. A sunset is beautiful because its visual harmony
increases coherence ($\gamma_{SE} \uparrow$, $\gamma_{AE} \uparrow$) at high integration.
Euler's identity $e^{i\pi} + 1 = 0$ is beautiful because it connects five previously
separate cognitive structures into one, increasing $P$ at $\Phi \gg 1$.

\paragraph{Meaning.}
$\mathrm{Meaning}_{\mathrm{peak}}(\Gam) = P \cdot D_{\mathrm{diff}} \cdot \Phi \cdot
R$~\status{I} --- the product of integrity, richness, connectedness, and self-awareness.
All four factors are necessary: if any equals zero, meaning vanishes. The multiplicative
form follows from the requirement that meaning is zero when any component is absent
(unlike a sum). Since meaning is a product, the most effective strategy for increasing it
is to boost the smallest limiting factor (Liebig's law applied to consciousness).
